% 数理逻辑（综述）
% license CCBYSA3
% type Wiki

本文根据 CC-BY-SA 协议转载翻译自维基百科\href{https://en.wikipedia.org/wiki/Mathematical_logic}{相关文章}。

数学逻辑（Mathematical Logic）是数学中对形式逻辑的研究。其主要分支包括模型论、证明论、集合论与递归论（又称可计算性理论）。数学逻辑的研究通常探讨逻辑形式系统的数学性质，例如其表达能力与推理能力。然而，它也可以用于刻画数学推理的正确性，或为数学建立理论基础。

自其诞生以来，数学逻辑既推动了数学基础研究的发展，也受到其驱动。该领域的起源可追溯至19世纪晚期，当时几何、算术与分析的公理化体系相继建立。20世纪初，大卫·希尔伯特（David Hilbert）提出了著名的“希尔伯特计划”（Hilbert’s Program），旨在证明数学基础理论的一致性。库尔特·哥德尔（Kurt Gödel）、格哈德·根岑（Gerhard Gentzen）等人的研究部分地解决了这一计划，同时也揭示了证明一致性问题的复杂性。

集合论的研究表明，几乎所有通常的数学内容都可以用集合语言形式化，尽管仍存在一些定理无法在常见集合论公理系统内被证明。现代数学基础研究的重点，往往不再是寻找能涵盖全部数学的理论，而是致力于分析哪些数学部分能够在特定的形式系统中被形式化（如“反向数学”所做的那样）。
\subsection{分支领域与研究范围}
1977年出版的《数学逻辑手册》（Handbook of Mathematical Logic由 J. Barwise 编辑，North-Holland 出版社，1977年\(^\text{[1]}\)）将当代数学逻辑大致划分为四个主要领域：
\begin{enumerate}
  \item 集合论（Set Theory）；
  \item 模型论（Model Theory）；
  \item 递归论（Recursion Theory）；
  \item 证明论与构造数学（Proof Theory and Constructive Mathematics），两者通常被视为同一研究方向的不同方面。
\end{enumerate}
此外，计算复杂性理论（Computational Complexity Theory）有时也被归入数学逻辑的研究范畴。\(^\text{[2][3]}\)
每个分支领域都有其独特的研究重点，但许多技术与结果在不同领域之间是共享的。各子领域之间的边界，以及数学逻辑与其他数学领域之间的界限，往往并不清晰。例如，哥德尔不完备定理不仅是递归论与证明论的重要成果，还引出了模态逻辑中的勒布定理（Löb's Theorem）；而强制法（Forcing）这一方法被广泛应用于集合论、模型论与递归论之中，并用于研究直觉主义数学。

范畴论（Category Theory）虽使用了大量形式公理化方法，并研究“范畴逻辑”（Categorical Logic），但通常不被视为数学逻辑的一个分支。由于范畴论在数学多个领域中具有广泛适用性，包括桑德斯·麦克莱恩（Saunders Mac Lane）在内的数学家提出将其作为独立于集合论的数学基础体系。这类基础体系依赖于“拓扑斯”（Topos）的概念，它可以被看作集合论的广义模型，并可在经典或非经典逻辑下定义。
\subsection{历史}
数学逻辑在19世纪中期作为数学的一个分支逐渐形成，源自哲学形式逻辑与数学传统的汇合。\(^\text{[4]}\)数学逻辑（又称“后勤学”（Logistic）、“符号逻辑”（Symbolic Logic）、“逻辑代数”（Algebra of Logic），或更近代的称呼“形式逻辑”（Formal Logic））是19世纪间发展起来的一系列逻辑理论的总称，它依赖一种人工符号系统与严格的演绎方法。\(^\text{[5]}\)在此之前，逻辑的研究主要与修辞学、计算术（calculationes\(^\text{[6]}\)）、三段论以及哲学研究相结合。20世纪上半叶，数学逻辑领域取得了爆炸性的发展，涌现出一系列奠基性成果，同时伴随着对数学基础问题的激烈争论。
\subsubsection{早期历史}
逻辑理论在历史上的许多文化中都得到了独立发展，包括古代中国、印度、希腊、罗马帝国以及伊斯兰世界。希腊的逻辑方法，尤其是亚里士多德的逻辑（即“项逻辑”，term logic），在其著作《工具论》（Organon）中系统呈现，对西方科学与数学的发展产生了长达数千年的深远影响。\(^\text{[7]}\)斯多葛派哲学家，尤其是克里西波斯（Chrysippus），开创了命题逻辑（propositional logic）的研究。18世纪的欧洲，一些哲学数学家——如莱布尼茨（Gottfried Wilhelm Leibniz）与兰伯特（Johann Heinrich Lambert）——尝试以符号化或代数化的方式处理形式逻辑的运算。然而，他们的研究成果大多零散，影响有限，并未引起广泛关注。
\subsubsection{十九世纪}
19世纪中叶，乔治·布尔（George Boole）与奥古斯都·德摩根（Augustus De Morgan）分别提出了系统化的数学逻辑理论。他们的研究基于代数学家乔治·皮科克（George Peacock）等人的工作，将传统的亚里士多德逻辑扩展为一个足以探讨数学基础的形式框架。\(^\text{[8]}\)1847年，瓦特罗斯拉夫·贝尔蒂奇（Vatroslav Bertić）独立于布尔，对逻辑的代数化也作出了重要贡献。\(^\text{[9]}\)此后，查尔斯·桑德斯·皮尔士（Charles Sanders Peirce）在布尔的成果基础上发展出一个包含关系与量词的逻辑体系，并于1870年至1885年间陆续发表了多篇论文。

哥特洛布·弗雷格（Gottlob Frege）则在1879年发表了《概念文字》（Begriffsschrift），独立提出了含量词的逻辑体系。这部著作通常被视为逻辑史上的分水岭。然而，弗雷格的工作在当时鲜为人知，直到20世纪初伯特兰·罗素（Bertrand Russell）开始推广其思想后才逐渐受到重视。弗雷格所创的二维逻辑符号体系虽然具有高度表达力，但并未被广泛采用，现代文献中几乎不再使用。

1890年至1905年间，恩斯特·施罗德（Ernst Schröder）出版了三卷本《逻辑代数讲义》（Vorlesungen über die Algebra der Logik）。  
该著作系统总结并扩展了布尔、德摩根与皮尔士的成果，是19世纪末符号逻辑研究的全面总结与重要参考文献。

\textbf{基础理论的发展}

数学未建立在坚实基础上的担忧，促使学者们为数学的基本领域（如算术、分析与几何）建立起严格的公理化体系。

在逻辑中，算术 一词通常指自然数理论。朱塞佩·皮亚诺（Giuseppe Peano）在1899年发表了著名的皮亚诺公理（Peano Axioms）\(^\text{[10]}\)，该体系在布尔与施罗德的逻辑系统基础上发展而来，并引入了量词（Quantifiers）。当时皮亚诺并不知道弗雷格的研究。与此同时，理查德·戴德金（Richard Dedekind）指出，自然数可以通过其数学归纳性质被唯一刻画。他提出了另一种非形式化的刻画方法，虽不具皮亚诺体系的逻辑严谨性，却能证明后者无法触及的若干定理，例如自然数集合的唯一性（同构意义下），以及从后继函数与归纳法定义加法和乘法的递归构造。\(^\text{[11]}\)

19世纪中叶，人们发现欧几里得几何体系的公理存在缺陷。\(^\text{[12]}\)除了尼古拉·罗巴切夫斯基（Nikolai Lobachevsky）于1826年证明的平行公设独立性\(^\text{[13]}\)外，数学家还发现某些欧几里得理所当然的命题实际上无法从其原有公理推导。例如，“一条直线至少包含两个点”，或“两圆若半径相等且圆心距离等于该半径，则必相交”。希尔伯特（David Hilbert）在帕施（Moritz Pasch）早期工作的基础上\(^\text{[15]}\)，提出了一套完整的几何公理体系\(^\text{[14]}\)。几何的成功公理化激发了希尔伯特进一步为数学的其他领域（如自然数与实数）寻求完全公理化描述的计划，这成为20世纪上半叶数学基础研究的重要方向。

19世纪的实分析理论取得了巨大进展，特别是在函数收敛性与傅里叶级数方面。  
卡尔·魏尔斯特拉斯（Karl Weierstrass）等数学家构造了“处处不可导但连续”的函数，挑战了人们对函数的传统理解（即函数是计算规则或光滑曲线）。他提出“分析的算术化”（arithmetization of analysis），试图以自然数的性质为基础公理化分析。 现代的\((\varepsilon, \delta)\)-极限定义最早由博尔查诺（Bernard Bolzano）于1817年提出\(^\text{[16]}\)，但鲜为人知。柯西（Augustin-Louis Cauchy）在1821年以无穷小定义连续性（见《分析教程》（Cours d’Analyse），第34页）。1858年，戴德金以有理数的“戴德金分割”（Dedekind cuts）定义实数，这一方法至今仍在教材中沿用。\(^\text{[17]}\)

格奥尔格·康托尔（Georg Cantor）发展了无限集合论的基本概念。他首先提出了基数理论，并证明了实数与自然数具有不同的基数。\(^\text{[18]}\)在随后的二十年间，康托尔发表了一系列论文，系统地发展了超穷数理论。1891年，他通过著名的“对角线论证”（diagonal argument）证明了实数集不可数，并进一步证明了“任意集合的幂集具有更大的基数”（即康托尔定理）。  康托尔相信“任何集合都可以良序化”，但始终未能证明该命题，遂于1895年将其列为公开问题。\(^\text{[19]}\)
\subsubsection{二十世纪的发展}
20世纪初，数学逻辑的主要研究方向集中在集合论与形式逻辑上。  
在非形式化集合论中发现的各种悖论，使一些学者开始怀疑数学体系本身是否存在内在的不一致性，从而促使他们寻求一致性证明。

1900年，希尔伯特（David Hilbert）在巴黎国际数学家大会上提出了著名的23个世纪问题，其中前两个分别要求：（1）解决连续统假设（Continuum Hypothesis）；  （2）证明初等算术的一致性。第十个问题则要求找出一种算法，能判定任意给定的多元整系数多项式方程是否存在整数解。这些问题的研究方向深刻影响了数学逻辑的发展。同样，希尔伯特于1928年提出的“判定问题”（Entscheidungsproblem）——即给定一个形式化的数学命题，是否存在一种算法能判定其真伪——也成为逻辑史上的核心议题。

\textbf{集合论与悖论的出现}

恩斯特·采梅洛（Ernst Zermelo）首先给出了“每个集合都可以良序化”的证明，这一结果此前连康托尔（Georg Cantor）也未能获得。\(^\text{[20]}\)为实现该证明，采梅洛引入了“选择公理”（Axiom of Choice），该公理一经提出便引发了数学界的激烈讨论与研究。针对其证明方法的质疑，采梅洛发表了第二篇论文，专门回应批评意见，并进一步澄清了证明过程。\(^\text{[21]}\)这篇论文促使选择公理逐渐被数学界广泛接受。

与此同时，关于选择公理的怀疑因“朴素集合论悖论”的发现而进一步加剧。切萨雷·布拉里-福尔蒂（Cesare Burali-Forti）首先提出了著名的布拉里-福尔蒂悖论（Burali-Forti Paradox），指出“所有序数的集合”本身无法构成集合。\(^\text{[22]}\)不久之后，伯特兰·罗素（Bertrand Russell）于1901年发现了罗素悖论（Russell’s Paradox），朱尔·理查（Jules Richard）也提出了理查悖论（Richard’s Paradox）。\(^\text{[23]}\)

采梅洛随后提出了首个集合论公理体系\(^\text{[24]}\)，后来亚伯拉罕·弗兰克尔（Abraham Fraenkel）加入了“替代公理”（Axiom of Replacement），二者共同形成了现代的采梅洛–弗兰克尔集合论（Zermelo–Fraenkel Set Theory，简称ZF）。采梅洛的体系通过“大小限制原理”（Limitation of Size Principle）有效避免了罗素悖论。

1910年，罗素与阿尔弗雷德·诺思·怀特海（Alfred North Whitehead）出版了《数学原理》（Principia Mathematica）第一卷。这部划时代的著作在类型论（Type Theory）的完全形式框架下建立了函数与基数理论，旨在彻底消除集合论悖论。尽管类型论未能成为数学的通用基础体系，但《数学原理》仍被视为20世纪最具影响力的著作之一。\(^\text{[25]}\)

随后，弗兰克尔26证明了：选择公理无法从采梅洛的集合论（带原子元素）中推出。  
保罗·科恩（Paul Cohen）\(^\text{[27]}\)后来进一步证明，即使不引入原子元素（urelements），选择公理在ZF体系中仍不可证。科恩在其证明中发展了“强制法”（Forcing），这一方法后来成为集合论中证明独立性命题的核心工具。\(^\text{[28]}\)

\textbf{符号逻辑}

利奥波德·勒文海姆（Leopold Löwenheim）\(^\text{[29]}\)与托拉夫·斯科伦（Thoralf Skolem）\(^\text{[30]}\)共同提出了著名的勒文海姆–斯科伦定理（Löwenheim–Skolem Theorem）。该定理指出：一阶逻辑（first-order logic）无法控制无限结构（infinite structures）的基数大小（cardinality）。斯科伦进一步意识到，这一结论同样适用于集合论的一阶形式化体系，并由此推导出：任何一阶集合论的形式化体系都必然存在一个可数模型（countable model）。这一违反直觉的结果后来被称为斯科伦悖论（Skolem’s Paradox）。
\begin{figure}[ht]
\centering
\includegraphics[width=6cm]{./figures/4c1b2e37ba280cee.png}
\caption{} \label{fig_SLlj_2}
\end{figure}
在其博士论文中，库尔特·哥德尔（Kurt Gödel）证明了完备性定理（Completeness Theorem），该定理建立了一阶逻辑（first-order logic）中语法（syntax）与语义（semantics）之间的对应关系。[31]哥德尔进一步利用完备性定理证明了紧致性定理（Compactness Theorem），从而揭示了一阶逻辑推理的有限性本质（finitary nature of logical consequence）。这些结果共同确立了一阶逻辑作为数学家使用的主导逻辑体系的地位。

1931年，哥德尔发表了划时代论文《论〈数学原理〉及相关体系中的形式上不可判定命题》（On Formally Undecidable Propositions of Principia Mathematica and Related Systems），其中证明了所有“足够强且可计算”（sufficiently strong and effective）的一阶理论的不完备性（incompleteness）。这一结论，即著名的哥德尔不完备定理（Gödel’s Incompleteness Theorem），揭示了数学公理体系的根本局限性，对希尔伯特（Hilbert）的“形式化一致性计划”（Hilbert’s Program）造成了沉重打击。  
它表明：在任何形式化的算术体系中，不可能在该体系内部证明其自身的一致性（consistency）。然而，希尔伯特起初并未充分认识到不完备定理的深远意义。[a]

哥德尔定理进一步表明：若一个形式体系是一致的（consistent），则其一致性无法在体系自身或任何更弱的体系中得到证明。这意味着——尽管理论体系的内部一致性不可被其自身形式化地验证——仍可能存在无法在体系内表达的外部一致性证明。此后，格哈德·根岑（Gerhard Gentzen）通过一个有限主义（finitistic）的系统，并辅以超穷归纳原理（transfinite induction），成功证明了算术体系的一致性。[32]根岑的研究引入了割消去法（cut elimination）与证明论序数（proof-theoretic ordinals）等核心概念，这些思想成为现代证明论（proof theory）的基础工具。  哥德尔则给出了另一种一致性证明方法：将经典算术的一致性归约为高阶直觉主义算术（intuitionistic arithmetic in higher types）的相容性问题。[33]

值得一提的是，第一本为普通读者撰写的符号逻辑教材早在1896年便已问世，其作者正是《爱丽丝梦游仙境》（Alice's Adventures in Wonderland）的作者路易斯·卡罗尔（Lewis Carroll）。[34][35]
